% Context from chapters-tex/variational-priors.tex; for mathematical reference, not compilation.
          
\subsection{Continuous Priors}
\label{sec-sob-tv-cont}

In this section we define priors for functions $f \in \Ldeux([0,1]^2)$ on a continuous domain. Section \ref{subsec-disc-priors} develops their counterparts for discrete images $f \in \RR^N$.

  
\paragraph{Sobolev prior.}

A prior energy assigns small values to images in the target class $\Theta$ and larger values, possibly $+\infty$, to images that depart from that model. The smoothness classes described in Section \ref{subsec-smooth-class} motivate Sobolev priors. The simplest such energy is
\eql{\label{eq-dfn-sob-energy}
	\Jsob(f)  = \frac{1}{2}|f|_{H^1}^2 = \frac{1}{2} \int \norm{\nabla f(x)}^2 \d x,
}
where $\nabla f$ is the distributional gradient and the integral is over $[0,1]^2$. The energy is one half of the squared $H^1$ seminorm and vanishes on constant functions. We set it to $+\infty$ when the gradient is not square integrable.

  
\paragraph{Total variation prior.}

To accommodate discontinuities in images, we use the total variation energy introduced in Section \ref{subsec-bv-model}. Rudin, Osher, and Fatemi introduced its use in image denoising \cite{rudin-tv}.

The total variation of a smooth image $f$ is defined as 
\eql{\label{eq-dfn-tv-energy}
	\Jtv(f) = \normTV{f} = \int \norm{\nabla_x f} \d x.
}
This energy extends to functions of bounded variation $f \in \text{BV}([0,1]^2)$, which may be discontinuous. This class includes indicator functions $f=1_{\Om}$ of sets $\Om$ with finite perimeter $|\partial \Om|$.

The total variation seminorm can also be computed using the coarea formula \eqref{eq-tv-coaera}, which shows in particular that $\normTV{1_{\Om}} = |\partial \Om|$.


                                                                      
\subsection{Discrete Priors}
\label{subsec-disc-priors}

An acquisition device discretizes a continuous-domain image $f \in \Ldeux([0,1]^2)$ into a pixel array $f \in \RR^N$. To process these data, we need priors defined on finite-dimensional image spaces.

  
\paragraph{Discrete gradient.}

Discrete Sobolev and total variation priors use finite-difference approximations to derivatives. For example, the forward differences are
\eqm{
	\de_1 f_{n_1,n_2} &= f_{n_1+1,n_2}- f_{n_1,n_2} \\
	\de_2 f_{n_1,n_2} &= f_{n_1,n_2+1}- f_{n_1,n_2},
}
where the image has $n\times n$ pixels and $N=n^2$. Higher-order schemes give more accurate approximations for smooth functions.
Boundary conditions require care. For simplicity, we use periodic boundary conditions, computing the indices $n_i+1$ modulo $n$. Symmetric boundary conditions can also be used to avoid boundary artifacts.

With the grid-spacing factor absorbed into the energy scaling, the discrete gradient at each pixel is
\eq{
	\nabla f_n = ( \de_1 f_n,\de_2 f_n ) \in \RR^{2}
}
and the gradient operator maps images to vector fields:
\eq{
	\nabla : \RR^N \longrightarrow \RR^{N \times 2}.
}
Figure \ref{fig-grad-discrete} shows discrete gradient vectors approximating the local direction of steepest intensity increase. Figure \ref{fig-discrete-operators} displays the gradient field and its magnitude. Large magnitudes indicate rapid intensity variations, typically at edges or in textured regions.

\myfigure{
\tabtrois{
\image{variational}{.3}{grad-discrete}&
\image{variational}{.3}{grad-discrete-zoom-1}&
\image{variational}{.3}{grad-discrete-zoom-2}
}
}{ 
	Discrete gradient vectors.
}{fig-grad-discrete}

\paragraph{Discrete divergence.}

One can also use backward differences,
\eqm{
	\tilde\de_1 f_{n_1,n_2} &= f_{n_1,n_2}- f_{n_1-1,n_2} \\
	\tilde\de_2 f_{n_1,n_2} &= f_{n_1,n_2}- f_{n_1,n_2-1}.
}
The adjoint relation between backward and forward differences is
\eq{
	\de_i^* = -\tilde \de_i,
} 
which means that
\eq{
	\foralls f,g \in \RR^N, \quad
	\dotp{\de_i f}{g} = - \dotp{f}{\tilde \de_i g},
}
This is the discrete counterpart of integration by parts,
\eq{
	\int_0^1 f' g = - \int_0^1 f g'
}
for smooth periodic functions on $[0,1]$.

We define the discrete divergence using backward differences:
\eq{
	\diverg(v)_n = \tilde \de_1 v_{1,n} + \tilde \de_2 v_{2,n},
}
This operator maps vector fields to images:
\eq{
	\diverg : \RR^{N \times 2} \longrightarrow \RR^{N}.
}
It is the negative adjoint of the gradient:
\eq{
	\diverg = - \nabla^*
}
which means that
\eq{
	\foralls f \in \RR^N, \;
	\foralls v \in \RR^{N \times 2}, \quad
	\dotp{\nabla f}{v}_{\RR^{N \times 2}} = -\dotp{f}{\diverg(v)}_{\RR^N}
}
This identity is a discrete version of the divergence theorem.

\myfigure{
\tabtrois{
\image{variational}{.3}{discrete-operators-image}&
\image{variational}{.3}{discrete-operators-grad}&
\image{variational}{.3}{discrete-operators-gradnorm}\\
Image $f$ & $\nabla f$ & $\norm{\nabla f}$
}
}{ 
	Discrete operators.
}{fig-discrete-operators}

  
\paragraph{Discrete Laplacian.}

The Laplacian is the divergence of the gradient:
\eq{
	\Delta f= \diverg( \nabla f ),
}
This operator is negative semidefinite because $\Delta=-\nabla^*\nabla$.

With the discrete gradient and divergence defined above, the Laplacian is the local high-pass filter
\eql{\label{eq-discrete-lapl}
	\Delta f_n = \sum_{ p \in V_4(n) } f_p - 4 f_n,
}
which, after scaling by the grid spacing, approximates the continuous Laplacian:
\eq{
	\pdd{f}{x_1}(x) + \pdd{f}{x_2}(x) \approx n^2 \Delta f_{k} \qforq x = k/n.
}

Laplacian operators thus act as filters. With the Fourier convention $\hat f(\omega)=\int f(x)e^{-\imath\langle x,\omega\rangle}\,\mathrm{d}x$, the continuous Laplacian is diagonal in the Fourier domain:
\eq{
	g = \Delta f \qarrq \hat g(\om) = -\norm{\om}^2 \hat f(\om)
}
while the discrete Laplacian \eqref{eq-discrete-lapl} has Fourier representation
\eql{\label{eq-discrete-lapl-fourier}
	g = \Delta f \qarrq \hat g_\om = -\rho_\om^2 \hat f_\om
	\qwhereq
	\rho_\om^2 = 4\sin\pa{ \frac{\pi}{n}\om_1 }^2 + 4\sin\pa{ \frac{\pi}{n}\om_2 }^2.
}

  
\paragraph{Discrete energies.}

The discrete Sobolev energy is one half of the squared $\ldeux$ norm of the gradient field:
\eql{\label{eq-discr-sobolev}
	\Jsob(f) = \frac{1}{2} \sum_n \big( ( \de_1 f_n )^2 + ( \de_2 f_n )^2 \big) = \frac{1}{2} \norm{\nabla f}^2.
}
The isotropic discrete TV energy instead sums the Euclidean magnitudes of the gradient vectors:
\eql{\label{eq-discr-tv}
	\Jtv(f) = \sum_n \sqrt{ ( \de_1 f_n )^2 + ( \de_2 f_n )^2 }
	= \normu{\nabla f}
}
Here the notation $\normu{v}$ for a vector field $v \in \RR^{N \times 2}$ means
\eql{\label{eq-lun-norm}
	\normu{v} = \sum_n \norm{v_n}
}
where $v_n \in \RR^2$.

                                                                      
                                                                      
                                                                      
\section{PDE and Energy Minimization}

Minimizing a smooth pr